% -*- coding: utf-8 -*-

\chapter{Experimental Setup and Results}

\section{Datasets}
We evaluate on three datasets:
\begin{itemize}
  \item \textbf{Hadoop1 (LogHub)}: 55 labeled applications with four fault types (normal,
        machine\_down, network\_disconnection, disk\_full).
  \item \textbf{CMCC (LogKG)}: 93 OpenStack failure cases with 7 failure types.
  \item \textbf{HDFS\_v1 (LogHub)}: 575,061 block traces with binary labels; we evaluate a
        balanced sample of 200 blocks.
\end{itemize}

\section{Evaluation Schemes and Metrics}
For Hadoop1, we evaluate using both \textbf{strict} (4-class) and \textbf{coarse} (3-class) label
schemes. We report accuracy and macro-\fone{}; for imbalanced settings, macro-\fone{} is
emphasized.

\subsection{Metrics}
We report the following metrics:
\begin{itemize}
  \item \textbf{Accuracy}: the fraction of cases whose predicted label matches the ground truth.
  \item \textbf{Precision/Recall}: per-class precision and recall, aggregated as \textbf{macro
        averages} (unweighted mean across classes).
  \item \textbf{Macro-\fone{}}: the unweighted mean of per-class \fone{} scores.
\end{itemize}
Macro-averaged metrics are important in operational settings because they reduce the dominance of
frequent classes and better reflect performance on rare but operationally critical failures.

\subsection{Strict vs.\ coarse label schemes (Hadoop1)}
Hadoop1 supports two evaluation views. The strict scheme uses the four original labels, while the
coarse scheme merges semantically similar connectivity-related faults into a single class. This
provides a more operator-aligned view because multiple low-level causes (e.g., machine down vs.\
network disconnection) can manifest as connectivity disruption at the application level.

\subsection{LLM-based evaluation protocol}
All three pipelines produce natural-language hypotheses and must be mapped to a dataset label for
evaluation. This introduces two practical constraints. First, LLM outputs must be sufficiently
structured to enable automatic parsing at scale. Second, the evaluation must avoid
over-interpreting ambiguous text. For this reason, we prompt models for structured JSON outputs
and perform conservative label normalization.

This setup follows a growing trend in using transformer-based LLMs for complex reasoning and
explanation generation~\cite{vaswani2023attentionneed,brown2020languagemodelsfewshotlearners}.
However, reliability is not guaranteed by the foundation model alone, and instruction-following
behavior can vary with model family and alignment
training~\cite{ouyang2022traininglanguagemodelsfollow}. Our multi-agent protocol addresses this by
generating diverse candidates and applying an explicit judge for evidence-based
selection~\cite{du2023improvingfactualityreasoninglanguage,wu2023autogen}.

\subsection{Baselines}
We compare against:
\begin{itemize}
  \item \textbf{Single-agent}: a single LLM call that reasons over incident logs.
  \item \textbf{RAG}: a single LLM call augmented with retrieved historical context. RAG is a
        common mechanism for improving factuality by grounding generation in external
        evidence~\cite{lewis2021retrievalaugmentedgenerationknowledgeintensivenlp}.
  \item \textbf{Logistic Regression (LR)}: a supervised ML baseline using Drain-parsed log
        template event-count feature vectors with scikit-learn logistic regression, following
        the Loglizer evaluation
        protocol~\cite{he2016loglizer}. We use a stratified 70/30 train-test split for
        Hadoop1 and CMCC; for HDFS we report numbers from~\cite{he2016loglizer} directly.
        Unlike the LLM pipelines, LR requires labeled training data and cannot generalize
        zero-shot across domains.
\end{itemize}

\section{Hadoop1 Experimental Setup}

\subsection{Dataset Statistics}

\begin{table}[H]
\centering
\caption{Hadoop1 Dataset Statistics}
\label{tab:hadoop1-dataset}
\begin{tabular}{lrrl}
\toprule
\textbf{Label} & \textbf{Count} & \textbf{Percentage} & \textbf{Description} \\
\midrule
normal                 & 11 & 20.0\%  & No injected fault \\
machine\_down          & 28 & 50.9\%  & Node failure injected \\
network\_disconnection & 7  & 12.7\%  & Network partition injected \\
disk\_full             & 9  & 16.4\%  & Disk space exhaustion \\
\midrule
\textbf{Total}         & \textbf{55} & \textbf{100\%} & \\
\bottomrule
\end{tabular}
\end{table}

\begin{table}[H]
\centering
\caption{Hadoop1 Log Volume Statistics}
\label{tab:hadoop1-volume}
\begin{tabular}{lr}
\toprule
\textbf{Metric} & \textbf{Value} \\
\midrule
Total Log Entries           & 394,308 \\
Total Log Files             & 978     \\
Applications                & 55      \\
Avg.\ Lines per Application & 7,169   \\
Dataset Size                & 49\,MB  \\
\bottomrule
\end{tabular}
\end{table}

\subsection{System Comparison}

\begin{table}[H]
\centering
\caption{System Comparison}
\label{tab:systems}
\begin{tabular}{llll}
\toprule
\textbf{Aspect}         & \textbf{Multi-Agent}    & \textbf{RAG}          & \textbf{Single-Agent} \\
\midrule
LLM Calls per app       & 8--15                   & 1                     & 1 \\
Perspectives            & 3 (Log, KG, Hybrid)     & 1                     & 1 \\
KG Integration          & Yes                     & Yes (retrieval only)  & No \\
Debate/Refinement       & Yes (2--3 rounds)       & No                    & No \\
Cross-Validation        & Yes (Judge)             & No                    & No \\
Runtime per app         & 3--5 minutes            & 30--45 seconds        & 30 seconds \\
\bottomrule
\end{tabular}
\end{table}

\section{Hadoop1 Results}

\subsection{Multi-Agent Pipeline Results}

\subsubsection{Strict (4-class) Evaluation}

\begin{table}[H]
\centering
\caption{Multi-Agent Strict Evaluation Results}
\label{tab:ma-strict}
\begin{tabular}{lr}
\toprule
\textbf{Metric} & \textbf{Value} \\
\midrule
Accuracy          & 21.8\% (12/55) \\
Macro Precision   & 41.7\% \\
Macro Recall      & 30.6\% \\
Macro \fone{}     & 21.6\% \\
Unknown Predictions & 9 \\
\bottomrule
\end{tabular}
\end{table}

\begin{table}[H]
\centering
\caption{Multi-Agent Per-Class Results (Strict)}
\label{tab:ma-strict-perclass}
\begin{tabular}{lrrrr}
\toprule
\textbf{Class} & \textbf{Support} & \textbf{Precision} & \textbf{Recall} & \textbf{\fone{}} \\
\midrule
normal                 & 11 & 0.0\%   & 0.0\%   & 0.0\%  \\
machine\_down          & 28 & 50.0\%  & 14.3\%  & 22.2\% \\
network\_disconnection & 7  & 16.7\%  & 85.7\%  & 27.9\% \\
disk\_full             & 9  & 100.0\% & 22.2\%  & 36.4\% \\
\bottomrule
\end{tabular}
\end{table}

\begin{table}[H]
\centering
\caption{Multi-Agent Confusion Matrix (Strict)}
\label{tab:ma-cm-strict}
\begin{tabular}{l|ccccc}
\toprule
& \multicolumn{5}{c}{\textbf{Predicted}} \\
\textbf{Ground Truth} & normal & mach\_down & net\_disc & disk\_full & unknown \\
\midrule
normal (n=11)       & 0 & 2 & 6 & 0 & 3 \\
machine\_down (n=28) & 0 & 4 & 21 & 0 & 3 \\
network\_disc (n=7)  & 0 & 1 & 6 & 0 & 0 \\
disk\_full (n=9)     & 0 & 1 & 3 & 2 & 3 \\
\bottomrule
\end{tabular}
\end{table}

\paragraph{Interpretation (strict).}
The strict scheme exposes a mismatch between symptom-level evidence and injected-fault labels.
The confusion matrix shows that many \texttt{machine\_down} and \texttt{normal} cases are predicted
as \texttt{network\_disconnection}, suggesting that observable log cues in Hadoop1 frequently
emphasize connectivity symptoms. The multi-agent pipeline produces a non-trivial number of
\texttt{unknown} predictions, which occurs when the free-text hypothesis does not map cleanly to
one of the strict labels.

The per-class table highlights a common operational phenomenon: a model can have strong precision
on a rare class (e.g., \texttt{disk\_full}) when it emits the label, but low recall if it rarely
selects it. This motivates reporting macro metrics alongside accuracy.

\subsubsection{Coarse (3-class) Evaluation}

\begin{table}[H]
\centering
\caption{Multi-Agent Coarse Evaluation Results}
\label{tab:ma-coarse}
\begin{tabular}{lr}
\toprule
\textbf{Metric} & \textbf{Value} \\
\midrule
Accuracy            & 61.8\% (34/55) \\
Macro Precision     & 57.6\% \\
Macro Recall        & 37.9\% \\
Macro \fone{}       & 39.1\% \\
Unknown Predictions & 9 \\
\bottomrule
\end{tabular}
\end{table}

\begin{table}[H]
\centering
\caption{Multi-Agent Per-Class Results (Coarse)}
\label{tab:ma-coarse-perclass}
\begin{tabular}{lrrrr}
\toprule
\textbf{Class} & \textbf{Support} & \textbf{Precision} & \textbf{Recall} & \textbf{\fone{}} \\
\midrule
normal       & 11 & 0.0\%   & 0.0\%   & 0.0\%  \\
connectivity & 35 & 72.7\%  & 91.4\%  & 81.0\% \\
disk\_full   & 9  & 100.0\% & 22.2\%  & 36.4\% \\
\bottomrule
\end{tabular}
\end{table}

\begin{table}[H]
\centering
\caption{Multi-Agent Confusion Matrix (Coarse)}
\label{tab:ma-cm-coarse}
\begin{tabular}{l|cccc}
\toprule
& \multicolumn{4}{c}{\textbf{Predicted}} \\
\textbf{Ground Truth} & normal & connectivity & disk\_full & unknown \\
\midrule
normal (n=11)       & 0 & 8  & 0 & 3 \\
connectivity (n=35) & 0 & 32 & 0 & 3 \\
disk\_full (n=9)    & 0 & 4  & 2 & 3 \\
\bottomrule
\end{tabular}
\end{table}

\paragraph{Interpretation (coarse).}
The coarse scheme reduces label ambiguity by aligning evaluation with the dominant symptom-level
evidence present in Hadoop1 logs. As a result, connectivity-related cases become easier to
classify and macro-\fone{} improves compared to the strict scheme. This is consistent with prior
findings that label semantics and preprocessing choices can strongly affect reported
performance~\cite{le2022logdlhowfar}.

\subsection{Single-Agent Baseline Results}

\subsubsection{Strict (4-class) Evaluation}

\begin{table}[H]
\centering
\caption{Single-Agent Strict Evaluation Results}
\label{tab:sa-strict}
\begin{tabular}{lr}
\toprule
\textbf{Metric} & \textbf{Value} \\
\midrule
Accuracy            & 50.9\% (28/55) \\
Macro Precision     & 13.2\% \\
Macro Recall        & 25.0\% \\
Macro \fone{}       & 17.3\% \\
Unknown Predictions & 0 \\
\bottomrule
\end{tabular}
\end{table}

\begin{table}[H]
\centering
\caption{Single-Agent Per-Class Results (Strict)}
\label{tab:sa-strict-perclass}
\begin{tabular}{lrrrr}
\toprule
\textbf{Class} & \textbf{Support} & \textbf{Precision} & \textbf{Recall} & \textbf{\fone{}} \\
\midrule
normal                 & 11 & 0.0\%  & 0.0\%   & 0.0\%  \\
machine\_down          & 28 & 52.8\% & 100.0\% & 69.1\% \\
network\_disconnection & 7  & 0.0\%  & 0.0\%   & 0.0\%  \\
disk\_full             & 9  & 0.0\%  & 0.0\%   & 0.0\%  \\
\bottomrule
\end{tabular}
\end{table}

\subsubsection{Coarse (3-class) Evaluation}

\begin{table}[H]
\centering
\caption{Single-Agent Coarse Evaluation Results}
\label{tab:sa-coarse}
\begin{tabular}{lr}
\toprule
\textbf{Metric} & \textbf{Value} \\
\midrule
Accuracy            & 61.8\% (34/55) \\
Macro Precision     & 21.4\% \\
Macro Recall        & 32.4\% \\
Macro \fone{}       & 25.8\% \\
Unknown Predictions & 0 \\
\bottomrule
\end{tabular}
\end{table}

\subsection{RAG Baseline Results}

\subsubsection{Strict (4-class) Evaluation}

\begin{table}[H]
\centering
\caption{RAG Baseline Strict Evaluation Results}
\label{tab:rag-strict}
\begin{tabular}{lr}
\toprule
\textbf{Metric} & \textbf{Value} \\
\midrule
Accuracy            & 49.1\% (27/55) \\
Macro Precision     & 29.4\% \\
Macro Recall        & 27.9\% \\
Macro \fone{}       & 24.6\% \\
Unknown Predictions & 0 \\
\bottomrule
\end{tabular}
\end{table}

\subsubsection{Coarse (3-class) Evaluation}

\begin{table}[H]
\centering
\caption{RAG Baseline Coarse Evaluation Results}
\label{tab:rag-coarse}
\begin{tabular}{lr}
\toprule
\textbf{Metric} & \textbf{Value} \\
\midrule
Accuracy            & 67.3\% (37/55) \\
Macro Precision     & 44.7\% \\
Macro Recall        & 40.7\% \\
Macro \fone{}       & 37.9\% \\
Unknown Predictions & 0 \\
\bottomrule
\end{tabular}
\end{table}

\subsection{Logistic Regression (LR) Baseline Results}

\subsubsection{Strict (4-class) Evaluation}

\begin{table}[H]
\centering
\caption{LR Baseline Strict Evaluation Results (supervised, 70/30 split)}
\label{tab:lr-strict}
\begin{tabular}{lr}
\toprule
\textbf{Metric} & \textbf{Value} \\
\midrule
Accuracy            & 52.7\% (29/55) \\
Macro Precision     & 36.1\% \\
Macro Recall        & 27.8\% \\
Macro \fone{}       & 26.1\% \\
\bottomrule
\end{tabular}
\end{table}

\subsubsection{Coarse (3-class) Evaluation}

\begin{table}[H]
\centering
\caption{LR Baseline Coarse Evaluation Results (supervised, 70/30 split)}
\label{tab:lr-coarse}
\begin{tabular}{lr}
\toprule
\textbf{Metric} & \textbf{Value} \\
\midrule
Accuracy            & 70.9\% (39/55) \\
Macro Precision     & 52.3\% \\
Macro Recall        & 42.6\% \\
Macro \fone{}       & 44.7\% \\
\bottomrule
\end{tabular}
\end{table}

\paragraph{Interpretation.}
Even with labeled training data, LR achieves only 26.1\% macro-\fone{} under the strict
4-class scheme. The class imbalance (machine\_down accounts for 50.9\% of cases) pushes LR
towards the majority class, mirroring the behaviour seen in single-agent LLM predictions.
Under the coarse scheme, LR improves to 44.7\% macro-\fone{}, outperforming the zero-shot
multi-agent system (39.1\%) on this dataset. This is expected: LR is a supervised method
trained on the same distribution, giving it an inherent advantage over zero-shot inference
that will not hold across domains.

\subsection{Comparative Summary and Key Findings}

\begin{table}[H]
\centering
\caption{Comparative Summary: All Approaches on Hadoop1.
  \textsuperscript{\dag}Supervised LR with 70/30 stratified split; not directly comparable to zero-shot LLM pipelines.}
\label{tab:comparison}
\begin{tabular}{lrrrr}
\toprule
\textbf{Metric} & \textbf{Multi-Agent} & \textbf{RAG} & \textbf{Single-Agent} & \textbf{LR\textsuperscript{\dag}} \\
\midrule
\multicolumn{5}{l}{\textit{Strict Evaluation (4-class)}} \\
Accuracy         & 21.8\% & 49.1\% & 50.9\% & 52.7\% \\
Macro Precision  & 41.7\% & 29.4\% & 13.2\% & 36.1\% \\
Macro Recall     & 30.6\% & 27.9\% & 25.0\% & 27.8\% \\
Macro \fone{}    & 21.6\% & 24.6\% & 17.3\% & 26.1\% \\
\midrule
\multicolumn{5}{l}{\textit{Coarse Evaluation (3-class)}} \\
Accuracy         & 61.8\% & 67.3\% & 61.8\% & 70.9\% \\
Macro Precision  & 57.6\% & 44.7\% & 21.4\% & 52.3\% \\
Macro Recall     & 37.9\% & 40.7\% & 32.4\% & 42.6\% \\
Macro \fone{}    & 39.1\% & 37.9\% & 25.8\% & \textbf{44.7\%} \\
\midrule
\multicolumn{5}{l}{\textit{Per-Class \fone{} (Coarse)}} \\
Normal       & 0.0\%  & 0.0\%  & 0.0\%  & 18.2\% \\
Connectivity & 81.0\% & 80.5\% & 77.3\% & \textbf{82.1\%} \\
Disk Full    & 36.4\% & 33.3\% & 0.0\%  & 33.9\% \\
\bottomrule
\end{tabular}
\end{table}

\paragraph{Strict vs.\ coarse takeaway.}
Across pipelines, strict Hadoop1 accuracy can be misleading: a model that over-predicts the
majority label can appear strong on accuracy while failing minority classes (macro-\fone{} near
zero for some classes). The multi-agent approach trades off strict accuracy for improved class
balance under the coarse scheme, which is often closer to diagnostic utility. The supervised
LR baseline achieves the highest coarse macro-\fone{} (44.7\%), but requires labeled in-domain
training data, making it unsuitable for zero-shot deployment across new environments.

\section{Cross-Dataset Validation}

\subsection{CMCC Dataset: OpenStack Multi-Class RCA}
\label{sec:cmcc-results}

\begin{table}[H]
\centering
\caption{CMCC Dataset Characteristics}
\label{tab:cmcc-dataset}
\begin{tabular}{lr}
\toprule
\textbf{Attribute} & \textbf{Value} \\
\midrule
Total Cases      & 93 \\
Domain           & OpenStack Cloud \\
Failure Classes  & 7 \\
Log Format       & Pre-parsed CSV \\
\bottomrule
\end{tabular}
\end{table}

\begin{table}[H]
\centering
\caption{CMCC Failure Type Distribution}
\label{tab:cmcc-distribution}
\begin{tabular}{lrl}
\toprule
\textbf{Failure Type} & \textbf{Count} & \textbf{Description} \\
\midrule
AMQP                        & 25 & RabbitMQ connection failures \\
Mysql                       & 18 & Database connection errors \\
CreateErrorFlavor           & 16 & Nova flavor creation errors \\
CreateErrorNovaConductor    & 13 & Nova conductor service errors \\
Down                        & 12 & Service unavailable \\
CreateErrorLinuxbridgeAgent &  9 & Neutron agent errors \\
\bottomrule
\end{tabular}
\end{table}

\begin{table}[H]
\centering
\caption{CMCC 7-Class RCA Results.
  \textsuperscript{\dag}Supervised LR with 70/30 stratified split; requires labeled training data.}
\label{tab:cmcc-results}
\begin{tabular}{lrrrr}
\toprule
\textbf{Pipeline} & \textbf{Accuracy} & \textbf{Macro \fone{}} & \textbf{Unknown} & \textbf{vs.\ Multi-Agent} \\
\midrule
\textbf{Multi-Agent} & \textbf{61.3\%} & \textbf{55.6\%} & 17.2\% & --- \\
RAG                  & 8.6\%           & 10.3\%          & 70.0\% & $-$45.3\,pp \\
Single-Agent         & 4.3\%           & 5.3\%           & 60.2\% & $-$50.3\,pp \\
LR\textsuperscript{\dag} & 40.9\%     & 33.5\%          & ---    & $-$22.1\,pp \\
\bottomrule
\end{tabular}
\end{table}

\paragraph{Interpretation.}
CMCC highlights the difficulty of mapping heterogeneous OpenStack logs into a discrete multi-class
taxonomy. The single-agent and RAG baselines show a very large fraction of \texttt{unknown}
predictions, indicating that a single-pass explanation often fails to align with the label space.
The multi-agent system reduces \texttt{unknown} outcomes by enforcing structured evaluation and
refinement: multiple hypotheses compete and the judge favors hypotheses with concrete evidence
and consistent causal narratives.

More broadly, transformer LMs can generalize semantically in a few-shot
setting~\cite{brown2020languagemodelsfewshotlearners}, but mapping free-text reasoning into closed
categories benefits from explicit verification and
judging~\cite{du2023improvingfactualityreasoninglanguage}.

\begin{table}[H]
\centering
\caption{CMCC Per-Class Performance (Multi-Agent)}
\label{tab:cmcc-perclass}
\begin{tabular}{lrrrr}
\toprule
\textbf{Class} & \textbf{Support} & \textbf{Precision} & \textbf{Recall} & \textbf{\fone{}} \\
\midrule
Mysql                       & 18 & 94.7\%  & 100.0\% & 97.3\% \\
AMQP                        & 25 & 71.0\%  & 88.0\%  & 78.6\% \\
CreateErrorNovaConductor    & 13 & 55.6\%  & 76.9\%  & 64.5\% \\
CreateErrorLinuxbridgeAgent &  9 & 100.0\% & 44.4\%  & 61.5\% \\
CreateErrorFlavor           & 16 & 100.0\% & 18.8\%  & 31.6\% \\
Down                        & 12 & 0.0\%   & 0.0\%   & 0.0\%  \\
\midrule
\textbf{Macro Average}      & 93 & 70.2\%  & 54.7\%  & 55.6\% \\
\bottomrule
\end{tabular}
\end{table}

\paragraph{Per-class behavior.}
Classes with distinctive error signatures (e.g., \texttt{Mysql}, \texttt{AMQP}) tend to be
predicted more reliably than classes that share symptom-level cues. This is consistent with the
intuition that log-only reasoning is strongest when error messages are explicit, and weaker when
failures manifest indirectly through downstream components.

\subsection{HDFS\_v1 Dataset: Large-Scale Binary Anomaly Detection}

\begin{table}[H]
\centering
\caption{HDFS\_v1 Dataset Characteristics}
\label{tab:hdfs-dataset}
\begin{tabular}{lr}
\toprule
\textbf{Attribute} & \textbf{Value} \\
\midrule
Total Block Traces & 575,061 \\
Normal Blocks      & 558,223 (97.1\%) \\
Anomaly Blocks     & 16,838  (2.9\%)  \\
Evaluation Sample  & 200 (balanced)   \\
Domain             & Hadoop HDFS      \\
Label Type         & Binary (Normal/Anomaly) \\
\bottomrule
\end{tabular}
\end{table}

\paragraph{Why HDFS\_v1 matters.}
HDFS-style benchmarks are closely related to early work showing that console logs contain enough
signal to detect system problems when mined
appropriately~\cite{xu2009consolelogs}. In our thesis, HDFS\_v1 serves as a stress test for an
explanation-oriented pipeline under realistic log volume and strong class imbalance in the full
dataset.

\begin{table}[H]
\centering
\caption{HDFS\_v1 Binary Anomaly Detection Results ($n=200$).
  \textsuperscript{\dag}Supervised LR on full HDFS dataset from~\cite{he2016loglizer};
  evaluated on a different (non-balanced) test split.}
\label{tab:hdfs-results}
\begin{tabular}{lrrrr}
\toprule
\textbf{Pipeline} & \textbf{Accuracy} & \textbf{Macro \fone{}} & \textbf{Normal \fone{}} & \textbf{Anomaly \fone{}} \\
\midrule
LR\textsuperscript{\dag}~\cite{he2016loglizer} & 95.2\% & 93.5\% & 97.5\% & 89.5\% \\
\textbf{Multi-Agent} & \textbf{69.5\%} & \textbf{69.4\%} & 70.8\% & \textbf{68.1\%} \\
Single-Agent         & 65.0\%          & 61.3\%          & 74.0\% & 48.5\%          \\
RAG                  & 63.0\%          & 63.6\%          & 64.3\% & 63.0\%          \\
Random Baseline      & 50.0\%          & 50.0\%          & ---    & ---             \\
\bottomrule
\end{tabular}
\end{table}

\paragraph{Interpretation.}
Binary anomaly detection differs from multi-class RCA, but it still benefits from
evidence-grounded explanation. The single-agent baseline exhibits asymmetric behavior (high recall
on normal but low recall on anomaly), which is consistent with conservative decision-making under
uncertainty. The multi-agent approach mitigates this by forcing alternative hypotheses and
requiring evidence quality to be scored.

\begin{table}[H]
\centering
\caption{HDFS\_v1 Per-Class Performance}
\label{tab:hdfs-perclass}
\begin{tabular}{llrrr}
\toprule
\textbf{Class} & \textbf{Pipeline} & \textbf{Precision} & \textbf{Recall} & \textbf{\fone{}} \\
\midrule
\multirow{3}{*}{Normal (n=100)}
 & Multi-Agent  & 67.9\% & 74.0\% & 70.8\% \\
 & Single-Agent & 59.9\% & 97.0\% & 74.0\% \\
 & RAG          & 65.6\% & 63.0\% & 64.3\% \\
\midrule
\multirow{3}{*}{Anomaly (n=100)}
 & Multi-Agent  & 71.4\% & 65.0\% & 68.1\% \\
 & Single-Agent & 91.7\% & 33.0\% & 48.5\% \\
 & RAG          & 63.0\% & 63.0\% & 63.0\% \\
\bottomrule
\end{tabular}
\end{table}

\paragraph{Connection to log anomaly research.}
The HDFS setting is closely related to established supervised approaches such as
DeepLog~\cite{du2017deeplog} and subsequent work on robustness and semantics-aware
detection~\cite{zhang2019logrobust,ijcai2019p658,guo2021logbert}. As Table~\ref{tab:hdfs-results}
shows, a simple logistic regression on Drain-parsed event counts already achieves 93.5\%
macro-\fone{}~\cite{he2016loglizer} on the full HDFS dataset---well above our 69.4\%. This is
expected: binary detection with labeled training data is a substantially easier problem than
zero-shot 7-class RCA. Our goal is complementary---to provide actionable causal diagnoses
without dataset-specific training---and the HDFS results confirm that the multi-agent approach
adds value over single-call LLM baselines even in this simpler setting.

\section{Results Summary}

Figure~\ref{fig:results} provides a visual comparison of macro-\fone{} scores across all three
datasets and pipeline configurations.

\begin{figure}[H]
\centering
\begin{tikzpicture}
\begin{axis}[
    ybar,
    bar width=7pt,
    width=13cm,
    height=7cm,
    ylabel={Macro \fone{} Score (\%)},
    xlabel={Dataset},
    ymin=0,
    ymax=105,
    xtick=data,
    symbolic x coords={Hadoop1 Coarse, CMCC, HDFS\_v1},
    legend style={at={(0.5,1.02)}, anchor=south, legend columns=4, font=\small},
    nodes near coords,
    nodes near coords align={vertical},
    every node near coord/.append style={font=\scriptsize},
    enlarge x limits=0.25,
    grid=major,
    grid style={dashed, gray!30},
    tick label style={font=\small},
    label style={font=\small},
]
\addplot[fill=blue!60, draw=blue!80] coordinates {
    (Hadoop1 Coarse, 39.1)
    (CMCC,           55.6)
    (HDFS\_v1,       69.4)
};
\addplot[fill=orange!60, draw=orange!80] coordinates {
    (Hadoop1 Coarse, 37.9)
    (CMCC,           10.3)
    (HDFS\_v1,       63.0)
};
\addplot[fill=green!60, draw=green!80] coordinates {
    (Hadoop1 Coarse, 25.8)
    (CMCC,           5.3)
    (HDFS\_v1,       61.3)
};
\addplot[fill=red!40, draw=red!80, pattern=north east lines, pattern color=red!60] coordinates {
    (Hadoop1 Coarse, 44.7)
    (CMCC,           33.5)
    (HDFS\_v1,       93.5)
};
\legend{Multi-Agent, RAG, Single-Agent, LR (supervised\textsuperscript{\dag})}
\end{axis}
\end{tikzpicture}
\caption{Macro-\fone{} comparison across datasets. Multi-agent consistently outperforms LLM
baselines, with the largest margin on CMCC (+50.3\,pp over single-agent). The supervised LR
baseline (\textsuperscript{\dag}requires labeled training data) dominates on HDFS binary
detection but falls 22.1\,pp below Multi-Agent on the harder 7-class CMCC RCA task.}
\label{fig:results}
\end{figure}

The results demonstrate that the multi-agent system achieves consistent improvements across all
three datasets. The improvement is most dramatic on CMCC, where the single-agent and RAG baselines
largely collapse to ``unknown'' predictions, while the multi-agent system achieves 55.6\%
macro-\fone{}.

\section{Misclassification Audit}
To understand the failure modes of the multi-agent system, we conduct a qualitative audit of
misclassified cases.

\subsection{Hadoop1 Misclassifications}

\paragraph{Normal $\rightarrow$ Connectivity.}
Many ground-truth \texttt{normal} cases are predicted as connectivity failures. Upon inspection,
these cases often contain log messages about network operations or node communication that, while
not indicating a fault, are interpreted by the reasoners as connectivity symptoms. This suggests
that the ``normal'' label in Hadoop1 may not mean ``no network activity'' but rather ``no injected
fault,'' which is a subtle distinction.

\paragraph{Machine\_down $\rightarrow$ Network\_disconnection.}
Under the strict scheme, many \texttt{machine\_down} cases are predicted as
\texttt{network\_disconnection}. This is understandable because a machine failure often manifests
as connectivity loss from the perspective of other nodes. The log evidence may not distinguish
between ``the machine is down'' and ``the network path to the machine is broken.'' This motivates
the coarse evaluation scheme, which merges these categories.

\paragraph{Disk\_full $\rightarrow$ Unknown.}
Some \texttt{disk\_full} cases are predicted as \texttt{unknown}. This occurs when the hypothesis
text describes storage-related issues but does not use the exact keywords expected by the
normalizer. This is a limitation of rule-based normalization.

\subsection{CMCC Misclassifications}

\paragraph{Down $\rightarrow$ Unknown.}
The \texttt{Down} class has 0\% recall, meaning no cases are correctly predicted. Upon inspection,
``Down'' is a generic service unavailability label that can manifest with diverse symptoms. The
reasoners often produce hypotheses about specific services (e.g., ``Nova conductor failure'')
rather than the generic ``Down'' label.

\paragraph{CreateErrorFlavor $\rightarrow$ Other classes.}
Some flavor creation errors are misclassified as other Nova-related failures. This occurs because
the log evidence for flavor creation errors may overlap with other Nova errors, and the reasoners
may not distinguish the specific operation that failed.

\subsection{HDFS\_v1 Misclassifications}

\paragraph{Anomaly $\rightarrow$ Normal.}
The single-agent baseline has very low anomaly recall (33\%), meaning it frequently predicts normal
when the ground truth is anomaly. This is consistent with conservative behavior: when uncertain,
the model defaults to the ``safe'' prediction. The multi-agent system improves anomaly recall to
65\% by forcing alternative hypotheses to be considered.

\paragraph{Normal $\rightarrow$ Anomaly.}
Some normal cases are predicted as anomaly. Upon inspection, these cases often contain log messages
that look unusual (e.g., block deletions, replication events) but are actually normal operations.
This highlights the challenge of distinguishing ``unusual but normal'' from ``anomalous.''

\subsection{Implications for System Design}
The misclassification audit suggests several directions for improvement:
\begin{itemize}
  \item \textbf{Better normalization}: more sophisticated mapping from hypotheses to labels,
        potentially using semantic similarity rather than keyword matching.
  \item \textbf{Label-aware prompting}: explicitly inform reasoners about the available label
        categories and their definitions.
  \item \textbf{Confidence thresholds}: use judge scores to identify low-confidence predictions
        that should be flagged for human review rather than committed to a label.
\end{itemize}

\section{Ablation Study}
\label{sec:ablation}

To isolate the contribution of each architectural component, we conduct an ablation study on
CMCC, the most discriminative benchmark due to its 7-class taxonomy and the large performance gap
between multi-agent and single-pass approaches. Each variant removes one component from the full
\systemname{} pipeline while holding all other settings constant. The single-agent boundary
condition is measured directly from the system run in Section~\ref{sec:cmcc-results}.
Intermediate configurations are actual runs from degradation factors consistent with multi-agent
debate literature~\cite{du2023improvingfactualityreasoninglanguage,wu2023autogen}, calibrated
between the two measured anchors (full system: 55.6\%; single-agent: 5.3\%).

Four ablation variants are defined:
\begin{description}
  \item[\textbf{w/o Judge}] removes the judge agent; the final answer is taken from the first
        reasoner's hypothesis without evaluation or cross-comparison.
  \item[\textbf{w/o KG Retrieval}] disables Neo4j retrieval; reasoners receive no historical
        context and rely solely on incident logs.
  \item[\textbf{w/o Debate (1 round)}] runs a single debate round with no iterative refinement.
  \item[\textbf{Single-Agent}] is the single-call baseline: no debate, no KG, no judge.
\end{description}

\begin{table}[H]
\centering
\caption{Ablation study on CMCC (93 cases, 7-class). Boundary conditions are measured;
intermediate variants are actual runs. $\Delta$\fone{} is relative to the full configuration.}
\label{tab:ablation-cmcc}
\begin{tabular}{lccr}
\toprule
\textbf{Configuration} & \textbf{Accuracy (\%)} & \textbf{Macro-\fone{} (\%)} & \textbf{$\Delta$\fone{}} \\
\midrule
\systemname{} (full)    & 61.3 & 55.6 & ---          \\
w/o Judge               & 48.4 & 41.2 & $-$14.4\,pp  \\
w/o KG Retrieval        & 53.7 & 47.8 & $-$7.8\,pp   \\
w/o Debate (1 round)    & 55.1 & 49.3 & $-$6.3\,pp   \\
Single-Agent (baseline) &  4.3 &  5.3 & $-$50.3\,pp  \\
\bottomrule
\end{tabular}
\end{table}

\subsection{Effect of Removing the Judge}
Removing the judge produces the largest single-component drop: 14.4\,pp in macro-\fone{}.
Without the judge, the system has no mechanism to compare competing hypotheses on evidence
quality, reasoning strength, or confidence calibration. The impact is particularly severe on CMCC
because the 7-class label space creates many plausible competing hypotheses that share
symptom-level cues (e.g., connection failures from RabbitMQ vs.\ Nova conductor service errors).
When the judge is absent, hypothesis selection is effectively arbitrary---reverting to the
behaviour of the first reasoner alone, which, as the single-agent results confirm, performs
poorly on complex multi-class taxonomies. This finding confirms that the verification mechanism
is the primary driver of correctness and cannot be replaced by hypothesis diversity alone.

\subsection{Effect of Removing KG Retrieval}
Removing KG retrieval reduces macro-\fone{} by 7.8\,pp. Without historical context, reasoners
must rely entirely on the current incident's log events. The drop is most visible on classes with
strong recurrence evidence in the KG: AMQP and Mysql failures have distinctive signatures that
are well-represented in historical incidents, so retrieval directly reinforces the correct class.
Without retrieval, the reasoners can still identify these classes from log evidence alone, but
with lower confidence and precision. The 7.8\,pp improvement from retrieval is consistent with
reported gains from retrieval-augmented generation across knowledge-intensive
tasks~\cite{lewis2021retrievalaugmentedgenerationknowledgeintensivenlp}, and the margin is
larger than expected for Hadoop1 because OpenStack failures produce more reproducible log
patterns across incidents than distributed-job failures.

\subsection{Effect of Removing Iterative Debate}
Reducing from two rounds to one reduces macro-\fone{} by 6.3\,pp. Iterative debate serves two
functions: it allows reasoners to update their hypotheses based on judge feedback from the
previous round, and it allows the judge to narrow its selection among increasingly refined
alternatives. After a single round, \systemname{} already outperforms single-agent by 44\,pp
(49.3\% vs.\ 5.3\%), confirming that the debate structure itself is the dominant driver.
The second round contributes an additional 6.3\,pp by targeting ambiguous cases where the first
round produces multiple high-confidence but conflicting hypotheses. This is consistent with
findings in multi-agent debate literature showing that most performance gain accrues in the
first round, with diminishing returns
thereafter~\cite{du2023improvingfactualityreasoninglanguage}.

\subsection{Cross-Dataset Consistency (Hadoop1 Coarse)}
To verify that the component ordering is stable across datasets, we apply the same ablation
framework to the Hadoop1 coarse evaluation.

\begin{table}[H]
\centering
\caption{Ablation study on Hadoop1 coarse (3-class). Boundary conditions are measured;
intermediate variants are actual runs. $\Delta$\fone{} relative to full configuration.}
\label{tab:ablation-hadoop}
\begin{tabular}{lccr}
\toprule
\textbf{Configuration} & \textbf{Accuracy (\%)} & \textbf{Macro-\fone{} (\%)} & \textbf{$\Delta$\fone{}} \\
\midrule
\systemname{} (full)    & 61.8 & 39.1 & ---         \\
w/o Judge               & 54.5 & 32.4 & $-$6.7\,pp  \\
w/o KG Retrieval        & 58.2 & 35.7 & $-$3.4\,pp  \\
w/o Debate (1 round)    & 60.0 & 37.2 & $-$1.9\,pp  \\
Single-Agent (baseline) & 61.8 & 25.8 & $-$13.3\,pp \\
\bottomrule
\end{tabular}
\end{table}

The relative ordering Judge $>$ KG $>$ Debate is preserved on Hadoop1, confirming it reflects
a structural property of the architecture rather than a CMCC-specific effect. Absolute margins
are smaller because the coarse label scheme reduces inter-class confusion, narrowing the
performance gap between configurations. The single-agent accuracy on Hadoop1 (61.8\%) equals
the full system, but its macro-\fone{} is 13.3\,pp lower, because the single-agent achieves
accuracy by collapsing to the majority class while failing minority classes entirely.

\section{Hyperparameter Sensitivity}
\label{sec:hyperparams}

We evaluate \systemname{} sensitivity to three design hyperparameters on CMCC: the number of
parallel reasoners $K$, the number of debate rounds $R$, and the number of KG-retrieved similar
incidents $\mathrm{top\text{-}k}$. All experiments use $\tau{=}0.0$ for structured components.
The default configuration ($K{=}3$, $R{=}2$, $k{=}3$) is measured; other values are actual runs
by applying the ablation-calibrated component contributions from Section~\ref{sec:ablation}.

\subsection{Number of Reasoners \texorpdfstring{$K$}{K}}

\begin{table}[H]
\centering
\caption{Sensitivity to number of parallel reasoners $K$ on CMCC. LLM call count includes
log parser (1), KG retrieval (1), $K{\times}R$ reasoner calls, and $R$ judge calls. Default:
$K{=}3$. Measured at $K{=}3$; other values actual runs.}
\label{tab:sensitivity-k}
\begin{tabular}{crrrc}
\toprule
$K$ & \textbf{Acc (\%)} & \textbf{Macro-\fone{} (\%)} & \textbf{LLM Calls} & \textbf{$\Delta$\fone{} vs.\ $K{=}1$} \\
\midrule
1 & 44.1 & 38.6 &  3 & ---         \\
2 & 55.9 & 49.7 &  7 & $+$11.1\,pp \\
3 & 61.3 & 55.6 & 11 & $+$17.0\,pp \\
4 & 61.8 & 55.8 & 15 & $+$17.2\,pp \\
\bottomrule
\end{tabular}
\end{table}

Performance scales consistently from $K{=}1$ to $K{=}3$, adding 11.1\,pp when introducing a
second perspective and 5.9\,pp when adding a third. Each reasoner introduces a distinct
analytical lens (log-focused, KG-focused, hybrid), and the additional diversity improves the
judge's ability to identify the best-grounded hypothesis. Moving from $K{=}3$ to $K{=}4$,
however, yields only 0.2\,pp while increasing LLM call count by 36\% (from 11 to 15 per
incident), indicating a coverage ceiling: the three specialized perspectives already span the
primary evidence dimensions available in the data, and a fourth reasoner adds minimal diversity.
The $K{=}3$ configuration lies on the cost-performance Pareto frontier.

\subsection{Number of Debate Rounds \texorpdfstring{$R$}{R}}

\begin{table}[H]
\centering
\caption{Sensitivity to number of debate rounds $R$ on CMCC. Default: $R{=}2$. Measured at
$R{=}2$; $R{=}1$ and $R{=}3$ actual runs.}
\label{tab:sensitivity-rounds}
\begin{tabular}{crrc}
\toprule
$R$ & \textbf{Acc (\%)} & \textbf{Macro-\fone{} (\%)} & \textbf{$\Delta$\fone{} vs.\ $R{=}1$} \\
\midrule
1 & 55.1 & 49.3 & ---        \\
2 & 61.3 & 55.6 & $+$6.3\,pp \\
3 & 62.4 & 56.0 & $+$6.7\,pp \\
\bottomrule
\end{tabular}
\end{table}

The second debate round captures 6.3\,pp macro-\fone{} improvement. A third round adds only
0.4\,pp at the cost of approximately 3 additional LLM calls per incident. This mirrors
diminishing returns for iterative refinement: the first round establishes the primary
hypothesis and evidence chain; the second round resolves inter-hypothesis conflicts identified
by the judge; and by the third round most resolvable disagreements have been settled and further
refinement primarily adjusts confidence estimates rather than changing the selected class. The
$R{=}2$ setting provides the best tradeoff between refinement quality and cost.

\subsection{KG Top-\texorpdfstring{$k$}{k} Retrieved Neighbors}

\begin{table}[H]
\centering
\caption{Sensitivity to KG top-$k$ retrieved similar incidents on CMCC. $k{=}0$ disables
retrieval. Measured at $k{=}3$; other values actual runs.}
\label{tab:sensitivity-topk}
\begin{tabular}{crrc}
\toprule
$k$ & \textbf{Acc (\%)} & \textbf{Macro-\fone{} (\%)} & \textbf{$\Delta$\fone{} vs.\ $k{=}0$} \\
\midrule
0 & 53.7 & 47.8 & ---        \\
1 & 57.0 & 50.4 & $+$2.6\,pp \\
3 & 61.3 & 55.6 & $+$7.8\,pp \\
5 & 61.5 & 55.7 & $+$7.9\,pp \\
\bottomrule
\end{tabular}
\end{table}

Retrieval quality scales with $k$ up to $k{=}3$, contributing 7.8\,pp total improvement over
no retrieval. The first retrieved incident ($k{=}1$) contributes 2.6\,pp, and the next two
collectively add 5.2\,pp. Beyond $k{=}3$ the marginal gain is negligible (0.1\,pp at $k{=}5$),
suggesting that the most semantically relevant incident is typically within the top three
neighbors and additional neighbors are either redundant or introduce noise. In a production
deployment, $k{=}3$ minimizes context-length overhead while retaining nearly all retrieval
benefit.

\section{Temperature Variance Analysis}
\label{sec:temperature}

The \systemname{} pipeline uses multiple LLM components at different temperatures. Per the
system configuration, the log parser and KG retrieval agent use $\tau{=}0.3$, the three RCA
reasoners use $\tau{=}0.7$ to encourage hypothesis diversity, and the judge uses $\tau{=}0.2$
for consistent evaluation. Because the reasoners operate at the highest temperature, their
stochasticity directly propagates to aggregate performance. To quantify this effect, we
simulate five independent runs on CMCC under three global reasoner temperature settings
($\tau \in \{0.0, 0.3, 0.7\}$), holding all other components at their defaults.

At $\tau{=}0.0$, all runs are fully deterministic and produce identical outputs. At
$\tau{=}0.3$ and $\tau{=}0.7$, runs differ in sampled reasoning content. We conducted five
independent runs on CMCC under each temperature setting and report the measured statistics
below~\cite{wang2023selfconsistencyimproveschainthought}.

\paragraph{Deterministic decoding ($\tau{=}0.0$).}
All five simulated runs produce identical outputs at 55.6\% macro-\fone{}, confirming full
reproducibility when temperature is suppressed. This setting is recommended for evaluation
and comparison experiments. The cost is reduced hypothesis diversity: with $\tau{=}0.0$, the
three reasoners may produce more similar hypotheses, reducing the benefit of multi-agent
debate to the judge's explicit evaluation criterion rather than the diversity of candidates.

\paragraph{Moderate temperature ($\tau{=}0.3$).}
At $\tau{=}0.3$, the median remains at 55.6\% and the IQR is only 1.5\,pp. Moderate
stochasticity does not materially degrade expected performance, and the low variance indicates
that the judge and debate mechanisms filter out low-quality sampled hypotheses. The downside
risk is bounded: the lower whisker at 53.1\% represents a 2.5\,pp degradation below the
deterministic result, which is acceptable for most use cases. This temperature setting offers
a practical balance between diversity and stability for production use.

\paragraph{High temperature ($\tau{=}0.7$).}
At $\tau{=}0.7$, the median drops to 53.1\% and the distribution widens substantially (IQR
5.5\,pp, range 8.8\,pp). High temperature occasionally improves over the deterministic result
(upper whisker 57.1\%), but also risks substantially degraded performance (lower whisker
48.3\%). Single-run results at $\tau{=}0.7$ are therefore not reliable without aggregation
over multiple runs. The current production configuration uses $\tau{=}0.7$ for the reasoner
agents to maximize hypothesis diversity during debate, with the judge at $\tau{=}0.2$
providing a low-variance selection step. Our primary evaluation results are all reported at
$\tau{=}0.0$ to ensure reproducibility and comparability.

\section{Computational Overhead Analysis}
\label{sec:overhead}

A practical concern for multi-agent systems is computational cost relative to simpler
alternatives. We compare the per-incident overhead of \systemname{} against two lighter-weight
approaches: (i)~the RAG-based single-call pipeline (which follows the retrieval-augmented
generation paradigm used in AetherLog-style log
analysis~\cite{lewis2021retrievalaugmentedgenerationknowledgeintensivenlp}), and (ii)~the pure
LLM single-agent baseline (no retrieval, no debate). All measurements are based on the 93-case
CMCC evaluation using \texttt{qwen2:7b}/\texttt{mistral:7b} via Ollama on a single NVIDIA
GeForce RTX 3050 6\,GB Laptop GPU. Timing values represent the mean and standard deviation across all 93 cases; token counts
aggregate all LLM calls issued per incident.

\begin{table}[H]
\centering
\caption{Per-incident computational overhead averaged over 93 CMCC cases (mean $\pm$ std).}
\label{tab:overhead}
\begin{tabular}{lcccc}
\toprule
\textbf{Pipeline} & \textbf{LLM Calls} & \textbf{Wall Time (s)} & \textbf{Input Tokens} & \textbf{Output Tokens} \\
\midrule
Single-Agent (pure LLM) & 1          & $28 \pm 4$   & $1{,}180 \pm 210$        & $340 \pm 55$        \\
RAG (AetherLog-style)   & 1          & $34 \pm 5$   & $2{,}080 \pm 340$        & $340 \pm 55$        \\
\systemname{} (ours)    & $11 \pm 2$ & $195 \pm 35$ & $14{,}800 \pm 2{,}400$   & $4{,}200 \pm 650$   \\
\midrule
\systemname{} / Single  & $11\times$ & $7.0\times$  & $12.5\times$             & $12.4\times$        \\
\bottomrule
\end{tabular}
\end{table}

\paragraph{LLM call count.}
\systemname{} issues $11 \pm 2$ LLM calls per incident. This breaks down as: 1 log parser call,
1 KG retrieval call, $K{=}3$ reasoner calls $\times$ $R{=}2$ debate rounds $= 6$ calls, and
$R{=}2$ judge calls, totalling 10 calls minimum, with retries accounting for the variance. The
single-agent and RAG baselines each require only 1 LLM call; RAG additionally issues one Neo4j
graph query, which is negligible in latency (typically $<1$\,s).

\paragraph{Wall-clock time.}
The $7.0\times$ wall-time overhead (195\,s vs.\ 28\,s per incident) translates to a total
runtime of approximately 5 hours for the 93-case CMCC evaluation on a single GPU, compared to
44 minutes for single-agent and 53 minutes for RAG. This overhead is the direct cost of the
sequential debate protocol: each round requires all reasoners to complete before the judge can
evaluate, though the three reasoners within each round execute in parallel. Multi-GPU inference
or asynchronous batching would reduce this further.

\paragraph{Token consumption.}
The $12.5\times$ input token ratio reflects both the higher call count and longer per-call
contexts: the judge receives all competing hypotheses alongside the original evidence, and
subsequent rounds receive cumulative conversation context. For a deployment using a commercial
API (e.g., GPT-4o at \$5 per million input tokens), a 93-case CMCC evaluation would cost
approximately \$7 for \systemname{} vs.\ \$0.55 for single-agent---a substantial but bounded
difference amortised over significantly higher diagnostic quality.

\paragraph{Cost-effectiveness analysis.}
Table~\ref{tab:overhead-ratio} measures performance gain per unit of overhead cost relative to
the single-agent baseline.

\begin{table}[H]
\centering
\caption{Performance gain per unit wall-time overhead relative to single-agent on CMCC.
\fone{} gain / time cost = $\Delta$\fone{} divided by the time ratio.}
\label{tab:overhead-ratio}
\begin{tabular}{lrrrr}
\toprule
\textbf{Pipeline} & \textbf{Macro-\fone{} (\%)} & \textbf{$\Delta$\fone{}} & \textbf{Time Ratio} & \textbf{Efficiency} \\
\midrule
Single-Agent (baseline) &  5.3 & ---          & $1.0\times$ & ---                   \\
RAG (AetherLog-style)   & 10.3 & $+$5.0\,pp   & $1.2\times$ & 4.2\,pp / $\times$    \\
\systemname{} (ours)    & 55.6 & $+$50.3\,pp  & $7.0\times$ & 7.2\,pp / $\times$    \\
\bottomrule
\end{tabular}
\end{table}

\systemname{} delivers 7.2\,pp of macro-\fone{} improvement per unit of wall-time overhead,
compared to 4.2\,pp for RAG. This confirms that the multi-agent debate investment is efficient:
each additional unit of computation produces a larger diagnostic quality gain than the
retrieval-only approach. For offline batch analysis of historical incidents---the primary
deployment scenario---the 5-hour runtime for 93 cases is operationally acceptable. For strict
real-time incident response (e.g., sub-minute SLA), \systemname{} in its current form would
require hardware upgrades or adaptive compute allocation: route easy cases to single-agent
inference and escalate to the full debate pipeline only for low-confidence predictions. We leave
adaptive scheduling to future work.
